The firmware is embedded C built with AVR-GCC and organised as a cooperative polling superloop with
no interrupt service routines. Each pass tests the keypad, services the active registration state
and tests the doctor's button. We chose this structure for auditability, since with no asynchronous
preemption the worst-case path through the loop can be reasoned about by inspection. The cost is
that every blocking operation, notably ADC acquisition and LCD refresh, extends the loop period for
all other activity, and that is what bounds the latency figures in Section~\ref{sec:results}.

\subsection{Registration State Machine}
The patient-side dialogue is a seven-state machine, shown in Fig.~\ref{fig:firmware-fsm}.
\textsf{IDLE} waits for a first keypress. The three data-entry states each own one field and one
prompt. The two acquisition states drive the LM35 and the pulse sensor and display a countdown so
the patient knows to hold still. \textsf{GENERATE\_SERIAL} allocates the serial number, enqueues the
record and returns to \textsf{IDLE}.

\begin{figure}[t]
\centering
\resizebox{\columnwidth}{!}{%
% Explicit coordinates throughout. Relative placement plus curved return paths
% was letting the "enqueue record" and "Prev" labels land on top of the IDLE
% box, so every node and label position below is pinned by hand.
\begin{tikzpicture}[
    st/.style={rectangle, draw, rounded corners, align=center, font=\small,
               minimum width=2.5cm, minimum height=0.66cm},
    lbl/.style={font=\scriptsize, inner sep=2pt},
    ar/.style={-{Latex[length=2mm]}, thick},
    dar/.style={-{Latex[length=2mm]}, thick, dashed}
]

% Forward path. Down the left column, across, then up the right column.
\node[st] (idle)  at (0,0)      {\textsf{IDLE}};
\node[st] (name)  at (0,-1.30)  {\textsf{ENTERING\_NAME}};
\node[st] (age)   at (0,-2.60)  {\textsf{ENTERING\_AGE}};
\node[st] (phone) at (0,-3.90)  {\textsf{ENTERING\_PHONE}};

\node[st] (temp)  at (5.10,-3.90) {\textsf{ENTERING\_TEMP}};
\node[st] (pulse) at (5.10,-2.60) {\textsf{ENTERING\_PULSE}};
\node[st] (gen)   at (5.10,-1.30) {\textsf{GENERATE\_SERIAL}};

\draw[ar] (idle)  -- node[lbl,right]{any key}      (name);
\draw[ar] (name)  -- node[lbl,right]{\textsf{Next}} (age);
\draw[ar] (age)   -- node[lbl,right]{\textsf{Next}} (phone);
\draw[ar] (phone) -- node[lbl,above]{\textsf{Next}} (temp);

% Right-column labels sit in the gap between the columns, clear of the
% horizontal Next arrow at y = -3.90.
\draw[ar] (temp)  -- node[lbl,left]{reading settled} (pulse);
\draw[ar] (pulse) -- node[lbl,left]{reading settled} (gen);

% Return path, routed well above the top row so nothing overlaps.
\draw[ar] (gen.north) -- (5.10,0.95) --
      node[lbl,above,pos=0.5]{enqueue record} (0,0.95) -- (idle.north);

% Prev corrections, drawn as short arcs to the left of the left column.
\draw[dar] (name.west)  to[out=170,in=190] (idle.west);
\draw[dar] (age.west)   to[out=170,in=190]
      node[lbl,left,pos=0.5]{\textsf{Prev}} (name.west);
\draw[dar] (phone.west) to[out=170,in=190] (age.west);

\end{tikzpicture}%
}
\caption{Patient-side registration state machine. Solid arrows are forward
  transitions and dashed arrows are \textsf{Prev} corrections.
  \textsf{GENERATE\_SERIAL} is the only state that enqueues, so a partial
  record can never enter the queue.}
\Description{A state diagram with seven boxes in two columns. The left column
reads from top to bottom: IDLE, ENTERING underscore NAME, ENTERING underscore
AGE, ENTERING underscore PHONE. An arrow labelled any key runs from IDLE down to
ENTERING NAME, and two further arrows labelled Next continue down the column. A
horizontal arrow labelled Next crosses from ENTERING PHONE to ENTERING
underscore TEMP at the foot of the right column. The right column then runs
upward, so an arrow labelled reading settled goes from ENTERING TEMP up to ENTERING
underscore PULSE, and a second reading settled arrow continues up to GENERATE
underscore SERIAL at the top of the right column. From there a long arrow
labelled enqueue record runs up, then left across the top of the figure, and
down into IDLE, closing the cycle. Three short dashed arcs on the far left curve
upward from ENTERING PHONE to ENTERING AGE, from ENTERING AGE to ENTERING NAME,
and from ENTERING NAME to IDLE; they are labelled Prev.}
\label{fig:firmware-fsm}
\end{figure}


Two properties follow from the topology, and both matter for correctness. Enqueueing happens on exactly
one transition, the one out of \textsf{GENERATE\_SERIAL}, which is reachable only after all five
preceding states complete, so a record abandoned halfway cannot reach the queue. That is the
structural guarantee behind the claim that partial records are impossible. Second, the
\textsf{Prev} key provides backward transitions among the data-entry states, letting a mistyped field
be corrected without restarting. The acquisition states offer no backward transition, because a
reading already written into the record is not editable by keypress.

\subsection{Alphanumeric Entry on a $4\times8$ Matrix}
Name entry is the one place a $4\times4$ keypad would have been inadequate, and the reason the design
pays twelve I/O lines for a $4\times8$ matrix. Scanning proceeds in two steps. With all eight columns
driven high, the row nibble identifies which row is asserted, after which each column is pulled low
in turn until the asserted row drops, identifying the column. The result is a key index of
$\mathrm{row} \times 8 + \mathrm{col}$.

Of the 32 keys, 26 carry the letters of the alphabet, four are editing and mode keys (space,
backspace, a caps-lock toggle and an alphabetic/numeric toggle), and two are the \textsf{Prev} and
\textsf{Next} navigation keys handled by the state machine rather than the character layer. The
mapping is direct rather than multi-tap. The first three rows cover A--X at one keypress each and the
fourth supplies Y and Z, so a name costs one press per character with no timing-based disambiguation
to learn. A three-character indicator (\texttt{abc}, \texttt{ABC} or \texttt{123}) in the LCD's
top-right corner keeps the active mode visible, and numeric mode remaps the first three columns to a
telephone-style digit layout. The age and mobile-number fields enter that mode automatically.

Contact bounce is handled in two stages. A detected press is re-tested after 20~ms and discarded if
it has not persisted, after which the scanner is locked out for 100~ms to suppress repeats from a
held key. Fields accumulate into a bounds-checked 16-character buffer, so an over-long entry
saturates rather than overruns.

\subsection{Queue Management}
A simple circular buffer was implemented in C to hold patient records. Each record contains the
assigned serial number; name, age and mobile number as entered at the keypad; the temperature
reading; and the pulse reading in BPM. The record is laid out as fixed-width character fields,
17~bytes each for name, age and mobile number, 10~bytes each for the two readings, and an integer
serial, for 73~bytes per patient. The buffer depth is fixed at compile time, so the static footprint
is $73 \times \mathrm{depth}$ bytes against the 2~KB of SRAM the ATmega32
provides~\cite{microchip2020}. Fixed-width fields make that footprint exactly predictable and
eliminate runtime allocation from the data path.

The enqueue operation is initiated following successful data entry and sensor acquisition. The
doctor's button initiates the dequeue process. Both execute in the same superloop on the same
structure, so there is exactly one writer and no synchronisation is required. The doctor-side control
has two-phase semantics that are not obvious from a single-button interface. The first press after
the queue becomes non-empty promotes the head record onto the console without removing it, so it
stays available for the duration of the consultation, and the next press retires it and promotes its
successor. One button therefore serves both ``show me this patient'' and ``I am done with this
patient''. The button is locked out for 500~ms after each press, and an emptied queue returns the
console to an explicit no-patient message rather than a stale record.

\subsection{Sensor Acquisition and Display}
Both sensors share the single ADC through the input multiplexer. For temperature the raw conversion
is scaled by the LM35 transfer function of 10~mV/$^{\circ}$C~\cite{ti2022} against the 2.56~V
reference, and successive samples are averaged to suppress conversion noise before the result is
committed. For heart rate the pulse waveform is sampled at fixed intervals and excursions above a
threshold are treated as beats, with $\text{BPM} = (60 \times 1000)/T_{\text{beat}}$ where
$T_{\text{beat}}$ is the interval in milliseconds between two consecutive detected peaks. The
acquisition state commits no value until the reading has settled, which is the firmware's response to
the motion sensitivity intrinsic to photoplethysmography~\cite{allen2007ppg}. Because the sampling
interval comes from calibrated delay loops rather than a crystal-driven timer, the time base carries
the $\pm 3\%$ tolerance of the internal RC oscillator, and Section~\ref{sec:results} treats this as
an error source in the reported accuracy.

The patient LCD carries the sequential prompts (``Enter Age'', ``Place Finger on Sensor'') followed by
the serial number, and the doctor LCD is driven over TWI using Fleury's I\textsuperscript{2}C master
library. Both use fixed field widths so a shorter value cannot leave residue from a previous patient
on screen, a real hazard on a display written in place rather than cleared. The firmware was flashed
over ISP with a USBasp programmer, which also set the CKSEL fuse bits, and the design was verified in
the Proteus model of Fig.~\ref{fig:proteus-complete} before assembly.
